\documentclass{amsart}

\usepackage{keytheorems}
\usepackage{amsthm}
\usepackage{color}

\newkeytheorem{problem}
\newtheorem*{proposition}{Proposition}
\newcommand{\todo}{\textcolor{red}{TO-DO.} }

\newcommand{\set}[1]{\left\{#1\right\}}
\newcommand{\powset}{{\cal P}}
\newcommand{\img}{_{*}}
\newcommand{\pimg}{^{*}}

\newcommand{\R}{\mathbb{R}}
\newcommand{\Z}{\mathbb{Z}}
\newcommand{\C}{\mathbb{C}}
\newcommand{\N}{\mathbb{N}}
\newcommand{\Q}{\mathbb{Q}}

\numberwithin{equation}{section}

\begin{document}

	\title{
		Math 251W: Foundations of Advanced Mathematics \\
		\textmd{Homework \#10: Portfolio Problems \S 5.2 \& 5.3}}
	\author{Teddy Wachtler}
	\date{\today}

	\maketitle

	\begin{problem}[manual-num=55]
		\begin{proposition}
			Let $n, q \in \N$, and let $a, b \in \Z$.
			Suppose $a \equiv b (\bmod{n}),$ and $q|n$.
			Prove that $a \equiv b (\bmod{q})$.
		\end{proposition}

		\begin{proof}
			Let $n$ and $q$ be arbitrary natural numbers.
			Let $a$ and $b$ be arbitrary integers.
			Suppose $a \equiv b (\bmod{n})$ and $q | n$.
			By the definition of modular congruence, we see that there exists a particular integer $c$ so that $a - b = cn$.
			By the definition of divides, we see that there exists a particular integer $d$ so that $n = qd$.
			By substitution, we see that $a - b = cqd$.
			By the closure of the integers, there exists a particular integer $e$ so that $e = cd$.
			By substitution, we see that $a - b = eq$.
			By the definition of modular congruence, we see that $a \equiv b (\bmod{q})$.
			Therefore, we see that for all natural numbers $n$ and $q$, and for all integers $a$ and $b$, if $a \equiv b(\bmod{n})$ and $q|n$, then $a \equiv b (\bmod{q})$.
		\end{proof}
	\end{problem}

\end{document}